% Replacement prose for Introduction and Related Work.
% Preserve the existing pipeline figure when integrating into main.tex.
% New citation keys are supplied separately in verified_additions.bib.
% Claims and numerical provenance: intro_related_中文对照与论证说明.md.

\section{Introduction}
Human videos offer a rich source of demonstrations for humanoid robots: a person can record a movement without specifying robot joint trajectories, collecting robot actions, or instrumenting a motion-capture studio. Turning these demonstrations into control requires connecting what a camera observes to what a different body can execute. Recent progress in world-grounded motion recovery~\cite{wham,gvhmr}, video-derived skill learning~\cite{videomimic}, and general motion tracking~\cite{beyondmimic,sonic} makes this connection increasingly practical. A central opportunity is to use the robot's own executed movements as visual training experience, so that human demonstrations and robot observations can contribute to the same video-to-control model.

Two developments shape this opportunity. Human-motion models provide a structured account of articulation and global movement, grounded in common body representations and large motion collections~\cite{smpl,amass,motionx}. Robot-motion methods translate this information into feasible references or learned control commands~\cite{gmr,nmr,sonic}. BeyondRetarget extends the mapping to monocular RGB, predicting robot trajectories through morphology-aligned supervision~\cite{beyondretarget}. These approaches establish effective routes from a demonstration to an executable target. They also motivate a complementary question: how can visual observations of the resulting robot execution become training inputs for the same interface that interprets a new human demonstration? This question concerns the organization of visual supervision around a reusable controller.

We introduce \method, a common visual motion interface learned from human videos and executable robot experience. The same frozen GVHMR estimator~\cite{gvhmr} processes human RGB and rendered robot RGB. Its estimated articulation, root movement, and stationary-joint cues form a common 145-dimensional motion representation; image features provide a learned residual for command prediction. A shared temporal mapper converts these observations into the skill commands of a frozen SONIC controller~\cite{sonic}. On robot images, the frontend produces a human-body motion proxy from pixels, making robot appearance pass through the same estimation path as human appearance. The controller combines the predicted commands with robot proprioception to generate joint targets. This design concentrates learning on the visual-to-command connection while reusing established motion perception and physical control.

Executable experience supplies the supervision for this connection. Offline, a human motion estimate enters the controller's existing motion interface; the resulting rollout records commands, actual robot states, and contacts. Rendering the executed movement produces a robot video paired with the original human source. Both observations can then train the shared mapper against execution-derived targets. Additional third-person videos can enter this process without synchronized motion capture for every source. Different renderings vary observation conditions, while additional source motions expand the demonstrated behaviors. At deployment, the learned mapper converts a new video into controller commands without a per-video explicit human-to-robot retargeting solve. The learning problem is therefore how to make executable visual experience useful across the two observation domains.

Our experiments examine this interface through physical execution, controlled frontend comparisons, and source-level scaling. On a fixed test set of 47 executable sources admitted from 96 candidates, increasing training from 256 to 1,411 source clips reduces human-input global robot joint-position error from 112.35 to 68.92\,mm at 120,000 updates, averaged over two seeds. This 38.7\% reduction is measured against execution-derived robot references. In a separate 12-clip locomotion development comparison, replacing the domain-specific visual frontend with the common motion frontend reduces robot-input joint-angle error by 19.0\% under the same training data and update schedule. Together, these results identify source coverage and the visual motion interface as useful design choices for learning video-conditioned control from executable experience.

Our contributions are:
\begin{itemize}
    \item A common visual motion interface that processes human and robot RGB with one frozen motion estimator and learns a shared temporal mapping to a reusable humanoid controller's skill commands, enabling new-video inference without an explicit human-to-robot retargeting solve.
    \item An executable paired-data construction process that connects human source videos with rendered observations, commands, states, and contacts from actual simulated robot execution, supporting supervision from expanded third-person video collections.
    \item A controlled empirical study of the interface and its data scale, demonstrating improved robot-input articulation with a common frontend and improved human-input execution accuracy as training source clips increase at fixed update budgets.
\end{itemize}

\section{Related Work}
\textbf{Human motion acquisition and representation.}
Motion acquisition provides the structured descriptions that connect visual observations with robot behavior. MoSh recovers body shape and motion from sparse markers~\cite{mosh}; SMPL supplies a learned parametric body model~\cite{smpl}; and AMASS unifies heterogeneous motion-capture collections through a common surface-based representation~\cite{amass}. These developments make motion data reusable across acquisition systems and learning tasks. Complementary datasets provide different forms of evidence: MoVi includes synchronized video and motion capture~\cite{movi}, Motion-X broadens expressive whole-body motion coverage~\cite{motionx}, and MotionPRO adds pressure observations to human motion capture~\cite{motionpro}. Their respective visual, geometric, and contact measurements support distinct aspects of motion estimation and evaluation.

Monocular reconstruction extends motion acquisition beyond instrumented capture. HMR regresses human shape and pose from a single image~\cite{hmr}, while VIBE uses temporal modeling and a motion discriminator to recover coherent sequences~\cite{vibe}. HMR~2.0 strengthens image-based mesh recovery with transformers and supports the tracking system of 4DHumans~\cite{hmr2}. Recovering a person's movement through the world additionally requires separating body and camera motion. SLAHMR optimizes human trajectories and camera motion together~\cite{slahmr}; WHAM combines image and temporal cues with contact-based trajectory refinement~\cite{wham}; and GVHMR estimates motion in coordinates defined by gravity and viewing direction~\cite{gvhmr}. We use GVHMR as a fixed visual motion frontend for both human and rendered robot videos. Our focus is the downstream use of its estimates as a common control input, including the human-body proxy it produces from robot appearance. This makes the choice of visual representation part of the control-learning problem.

\textbf{Human-to-robot motion mapping.}
Transferring motion between bodies requires accounting for morphology, joint limits, and physical feasibility. Early humanoid systems adapted human gestures to robot motion ranges~\cite{pollard2002} and generated balanced whole-body movements from captured dances~\cite{nakaoka2003}. Ayusawa and Yoshida jointly optimized morphing parameters and motion under robot constraints~\cite{ayusawa2017}. More recently, GMR combines motion scaling with staged optimization and studies how reference quality affects downstream tracking~\cite{gmr}. NMR learns a human-to-robot mapping from trajectories produced through geometric retargeting and physics-based refinement~\cite{nmr}. These methods establish explicit robot-motion references as an effective interface between human demonstrations and control, with either optimization or a learned model performing the conversion.

Learned interfaces offer several output choices. SONIC encodes human, robot, and hybrid motion commands into shared tokens for a state-conditioned whole-body controller~\cite{sonic}. Its human-motion input already supports deployment without an external per-frame retargeting solve. BeyondRetarget predicts robot root motion and joint configurations directly from monocular video, using morphology-aligned supervision and contact refinement before execution with a motion tracker~\cite{beyondretarget}. The official Plus implementation instead combines a frozen human-motion estimator with robot adapters and global-trajectory refinement~\cite{beyondretargetpluscode}. Our work builds on the established SONIC command interface and learns its visual input mapping from human and executed robot observations. The distinguishing design is that both appearances pass through the same frozen motion estimator and train the same command mapper. Robot executions thus supply visual inputs as well as control targets, connecting motion conversion with learning from paired visual experience.

\textbf{Video-based imitation and physical control.}
Visual imitation has long studied representations that connect human demonstrations with robot observations. Time-Contrastive Networks learn from synchronized views and use visual representations for pose imitation and reward learning~\cite{tcn}. XIRL learns task-progress representations across embodiments and converts them into rewards for policy learning~\cite{xirl}. These works establish the value of visual correspondence across bodies. We investigate a more specific interface for humanoid motion: a fixed motion estimator supplies structured inputs, execution supplies paired command supervision, and a fixed skill controller supplies the motor capabilities.

Recent humanoid systems also differ in the role of video. VideoMimic reconstructs human movement and scene geometry from demonstrations, then learns contextual skills that respond to the robot's environment~\cite{videomimic}. RoboMirror uses video-derived motion intent to condition a diffusion policy for humanoid locomotion~\cite{robomirror}. Our target is the motion in a supplied demonstration, expressed through a reusable controller. That controller belongs to a broader line of physical imitation: DeepMimic learns dynamic skills from example motions~\cite{deepmimic}, BeyondMimic connects motion tracking with guided motion generation~\cite{beyondmimic}, and SuperMimic trains whole-body control through a differentiable world model~\cite{supermimic}. Holding the controller fixed allows our study to measure how visual processing, paired robot observations, and training-source coverage affect the executed movement. This complements advances in the controller's own tracking and skill repertoire.
